\documentclass[11pt,a4paper]{article}
\usepackage[T1]{fontenc}
\usepackage[utf8]{inputenc}
\usepackage[margin=1.25cm]{geometry}
\usepackage{enumitem}
\usepackage{amsmath,amssymb}

\begin{document}

\begin{center}
{\Large\bfseries Answer Key}\\[2pt]
{\large Time Travellers! --- Calendars, conversions and timelines}
\end{center}

\section*{Part A: Calendar hops}

\begin{enumerate}[label=\textbf{A\arabic*.}, leftmargin=*, itemsep=2pt]
\item Saturday 7 November 2026.
\item Tuesday 10 November 2026.
\item Saturday 17 October 2026.
\item Monday 23 November 2026.
\item Wednesday 16 September 2026.
\item Wednesday 4 November 2026 at 14:00.
\item 50 days \(=\) 7 weeks 1 day.
\item No. Nearest reachable Saturdays: 19 December 2026 and 26 December 2026.
\item Friday 13 November 2026.
\end{enumerate}

\section*{Part B: Conversion station}

\textbf{1. Conversions.}
\begin{itemize}[leftmargin=*, itemsep=1pt]
\item \(3\) weeks \(=21\) days.
\item \(4\) centuries \(=40\) decades.
\item \(2\) days \(=48\) hours.
\item \(2\) millennia \(=20\) centuries.
\item \(180\) minutes \(=3\) hours.
\item \(7200\) seconds \(=2\) hours.
\item \(5\) decades \(=50\) years.
\item \(3\) days \(=4320\) minutes.
\end{itemize}

\textbf{2. Dial readings.}
\begin{enumerate}[label=\textbf{B\arabic*.}, leftmargin=*, itemsep=2pt]
\item 3 days.
\item 60 weeks.
\item 41 days 16 hours.
\item 11 days 13 hours 46 minutes 40 seconds.
\end{enumerate}

\textbf{3. Leap years.} 2028 and 2400.

\textbf{4. Order from shortest to longest.}
\[
1{,}000{,}000\text{ s} < 50{,}000\text{ min} < 0.1\text{ years} < 1000\text{ hours} < 6\text{ weeks}.
\]

\section*{Part C: Timeline jumps}

\textbf{1. Big timeline marks.} Tutankhamun: 1323 BC (between 1500 BC and 1000 BC, closer to 1500 BC); First Olympic Games: 776 BC (between 1000 BC and 500 BC); Roman invasion: AD 43 (just right of the 1 BC/AD 1 mark); Battle of Hastings: AD 1066 (just after AD 1000).

\textbf{2. Years ago from 2026.} (a) 960 \quad (b) 1983 \quad (c) 2801 \quad (d) 4585.

\textbf{3. 2(a) in centuries and decades.} 9 centuries and 6 decades.

\textbf{4. Landing places.} 975 BC; AD 2526; 475 BC.

\textbf{5. Zoomed timeline marks.} Crosses at 1215, 1666, 1969, and 2026.

\textbf{6. Differences.} (a) 451 years \quad (b) 57 years.

\textbf{7. Closer to today.} The Great Fire is closer to today: 360 years versus 600 years to Hastings.

\section*{Part D: Mission log}

\textbf{Mission table.}
\begin{itemize}[leftmargin=*, itemsep=1pt]
\item Jump 1: BACK 3 centuries, AD 2026 \(\to\) AD 1726; battery used \(6\%\).
\item Jump 2: FORWARD 25 years, AD 1726 \(\to\) AD 1751; battery used \(0.5\%\).
\item Jump 3: BACK 2 millennia, AD 1751 \(\to\) 250 BC; battery used \(40\%\).
\item Jump 4: FORWARD 1500 years, 250 BC \(\to\) AD 1251; battery used \(30\%\).
\end{itemize}

\begin{enumerate}[label=\textbf{\arabic*.}, leftmargin=*, itemsep=2pt]
\item FORWARD 775 years.
\item \(3825\) years \(=38.25\) centuries.
\item \(76.5\%\) used; \(23.5\%\) left.
\item Yes. Needs \(15.5\%\); arrives with \(8\%\) left.
\item One possible plan:
\begin{center}
\begin{tabular}{c|c|c|c}
Jump & Dial setting & Lands in & Battery used\\
\hline
1 & BACK \(4585\) years & 2560 BC & \(91.7\%\)\\
2 & FORWARD \(3625\) years & AD 1066 & \(72.5\%\)\\
3 & FORWARD \(903\) years & AD 1969 & \(18.06\%\)\\
Home & FORWARD \(57\) years & AD 2026 & \(1.14\%\)\\
\hline
\multicolumn{3}{r}{Total} & \(183.4\%\)\\
\end{tabular}
\end{center}
Total battery needed is \(183.4\%\), so a recharge is required. Any order visiting all three and returning needs at least \(9170\) years of travel.
\end{enumerate}

\end{document}